\subsection{三高分母の完全マッチング定理：今後は仮定でなく入力として使う}
\label{nm18:matching}

整数行和・pair products を満たし，$q_0=1,q_1,q_2,q_3>1$ である局所模型を考える．
正の entries の無向 graph は，完全マッチングまたは低 ray を含む三角形と孤立頂点の和に限られる．
さらに，正の basin weights を持つことと，完全マッチング型であることは同値である．
分母上限，clutching，node，対向 $D$ 条件はいずれも不要である．
証明は算術稿の \texttt{mh18:matching} に置く．

再利用時には，低 row の三つの entries がちょうど二負一正であることを最初に課す．
正の低 entry の相手を $c$，負の二つの相手を $a,b$ とすると，
正符号 graph は $0c\sqcup ab$ であり，既存の canonical Schur 表示が全ての正の三高模型を覆う．
$\beta\le-1$ と，境界 $\beta=-1$ の完全な約数表示をそのまま使える．
ただし，全ての高 $D$ が非負となる可能性はまだ解決していない．
この係数問題の残りは $\beta\le-2$ であり，特異な Schur 系も保持する．

\paragraph{他の問題にも使える二つの証明操作．}
二つの道型では，同じ残余二次式
\[
 (kC-b)L^2+\{kd-aC+B(kC-b)\}L+d(kB-a)=0
\]
に対し，全係数の符号がそろう．正の根 $L$ がないことが，無上界の排除を与える．
三角形型では，正の三閉路の不等式を
\[
 (tS-H)K<1,\qquad (sT-K)H<1
\]
と書くと，整数性から $H=tS+\alpha,K=sT+\beta$，$\alpha,\beta\ge0$ となる．
残る pair equations を整理した
\[
 bS=AT+\alpha\beta,\qquad
 (A-1)U=CS+k(\alpha+1),
\]
\[
 (b+1)U=(C+1)T+(k+1)(\beta+1)
\]
は $A>1,b>0,U>T$ を強制する．$\delta=U-T\ge1$ に対して
\[
 A\{s(1-\delta)-k\beta\}=bsU+bk(s+t)+Ct\beta
\]
の左辺は非正，右辺は正となる．
ここでは $A=s\alpha-t$，$b=s-t\beta$，$C=s+(s+t+1)k$ であり，
この排除に用いる変数は canonical matching 型の変数と混同しない．
実数での $U>T$ を整数の $\delta\ge1$ へ強める一歩が不可欠である．

\paragraph{正値性を省く誤りを避ける検査例．}
分母 $(1,5,19,31)$ と
\[
 P=\begin{pmatrix}
 0&1&1&-3\\4&0&2&-11\\18&47&0&-84\\-10&-14&-7&0
 \end{pmatrix}
\]
は整数行和・pair products を全て満たし，正符号 graph は三角形である．
しかし閉路 $0\to2\to1\to0$ の gain は $188/95>1$ であり，正 weights はない．
行和・pair products だけから完全マッチングを結論してはいけない．

\nolinkurl{verify_three_high_matching_v18.py} は十一の多項式恒等式を係数ごとに確認し，
道型の係数の符号を正変数の平行移動で検査する．
併せて固定低データの小さい入力を既存の完全復元器へ通し，実装の整合性を点検する．
この入力格子の検査を，一般定理の証明として引用しない．
